The structural-mechanics experiments used CPU-only cloud instances,
with single-precision network training and double-precision kernel solves
and error computations. The second experiment of Section~\ref{sec:seeds} ran
on one Intel Xeon E5-2698 v4 host (20 physical cores, 40 hardware threads).
The distributed \texttt{StructuralMechanics} dataset contains 40000
samples. Its inputs have an exact broadcast structure and are represented
in $\mathbb R^{41}$. Samples 1--20000 form the training block and samples
20001--40000 form the test block. In the high-data protocol, a fixed
permutation assigns 1000 training-block samples to validation and 19000
to fitting. The low-data protocol uses the first 1250 training-block
samples, with the final 250 reserved for validation. Model selection uses
the validation split.

\paragraph{FNO.} The network has width 64, 14 modes per dimension, and
4 spectral layers with pointwise linear skips and GELU activations. The
input lift takes 3 channels (the broadcast load and two coordinates), and
the output projection has dimensions $64\to128\to1$, giving 6.45M
parameters in total. Training uses AdamW with learning rate
$1.5\times10^{-3}$, weight decay $10^{-6}$, a cosine schedule ending at
$10^{-6}$, and batch size 128. The first experiment scheduled 200 epochs
and retained the best validation checkpoint reached before its time limit;
the second completed a 100-epoch schedule.

\paragraph{Transformer.} This architecture is excluded from the reported
ensemble. Its encoder uses 41 tokens, each comprising the load value and
10 sine/cosine pairs of the coordinate, with 5 pre-norm blocks, dimension
192, 4 attention heads, and MLP ratio 4. The decoder forms 1681 grid
queries from Fourier features of both coordinates and uses two
cross-attention blocks followed by a pointwise $192\to192\to1$ head.
The network has 3.16M parameters.

\paragraph{UNet.} The network has three convolutional scales
($41\to21\to11$ by average pooling with ceiling mode) and a $6\times6$
bottleneck. Each block uses two $3\times3$ convolutions with GroupNorm(8)
and SiLU, with widths $48/96/192/384$. Bilinear upsampling, skip
concatenation, and a $1\times1$ output head complete the network, which
has 4.38M parameters and the same input channels as the FNO. Training uses
AdamW with learning rate $1.5\times10^{-3}$, weight decay $10^{-5}$, a
cosine schedule ending at $10^{-6}$, and batch size 256. The second
experiment uses 100 epochs.

\paragraph{MSE variant.} The residual MLP uses mean squared error on
standardized targets in place of the metric loss, with a 120-epoch schedule
and otherwise the same training settings. It has the lowest error among the plain MLPs and introduces a second
training loss into the ensemble.

\paragraph{MLP.} The input layer has dimensions $41\to1024$, followed by
three residual SiLU blocks of width 1024 and an output layer
$1024\to1681$, for 4.91M parameters. Training uses AdamW with learning
rate $10^{-3}$, weight decay $10^{-5}$, a cosine schedule, 400 epochs, and
batch size 256. A wider variant has width 1536, five residual blocks and
14.5M parameters, with the same training protocol. The low-data pipeline
uses the width-1024 MLP for 400 epochs with batch size 64 on the fixed
1000 fitting rows.

\paragraph{Refiner.} The load is concatenated with its kernel stress
prediction and passed through an input layer of dimensions
$(41+1681)\to1024$ and the same residual trunk as the MLP. The resulting
network has 6.64M parameters and uses the MLP training settings for 100
epochs. Its conditioning channel consists of four-fold out-of-fold kernel
predictions during training and full-data kernel predictions at evaluation.
Reflection
augmentation flips the load and permutes both the kernel channel and the target
by the row-reversal index. Evaluation averages $F(u,h(u))$ with
$T F(Su,T h(u))$; it does not recompute the kernel prediction at $Su$.
The full-map guarantee and the additional
condition on the kernel channel are distinguished in
Supplement~\ref{sec:theory-sym}.

\paragraph{Common training details.} Except for the MSE variant, the loss
is the batch mean of $\|\hat v-v\|_2/\|v\|_2$ in original units. Targets
are standardized by the per-pixel training mean and global standard
deviation, and predictions are denormalized when evaluating the loss.
Neural members use reflection augmentation with probability $1/2$ and
reflection averaging at evaluation, with the supplied-field convention
for the refiner described above. The kernel-only member uses neither
operation. Validation error is evaluated every 10 epochs and the checkpoint
with the lowest error is retained. Section~\ref{sec:seeds-arch} compares
residual alignment within and between configurations in the sixty-predictor
pool.

\paragraph{Test-block use.} The structural-mechanics test block is the
benchmark's fixed public split. It is used by the first structural-mechanics experiment, the
six-member experiment, the ablations, the hard-case study, the learning curves,
the sixty-predictor analysis, the conformal calibration (which draws its
$1000$ calibration cases from it) and the additional controls. Predictor
fitting and model-level selection use the fitting and validation splits.
The test-block analyses additionally fit the hindsight ensemble optimum and choose calibration quantiles on a test subset.
The same $2000$ public OCO-2 test states are reused across the ten-seed
experiment, the readout control and the ensemble comparison.

\paragraph{Training labels and target centering.}
The pooled-centering runs use \texttt{krr\_oof.py} with
\texttt{--target-centering} \texttt{pooled}; the fold-local variant uses
\texttt{--target-centering} \texttt{fold-local}. Only the kernel channel uses foldwise fitted kernel weights. Its predictions use a pooled target mean over all 19000 fitting rows, so held-out fold targets enter that centering statistic. Neural training predictions are generated by the networks fitted on those rows, and the refiner uses its training kernel channel. The correction is therefore fitted to $Y-M_{\rm tr}$. At query time the mean uses full-data kernel predictions. Section~\ref{sec:theory-or} gives the explicit $m_{\rm full}(X)-M_{\rm tr}$ mismatch term. Supplement~\ref{supp:sensitivity} measures its propagation with the fitted mean and kernel fixed, and compares target-centering choices after repeating refiner training and downstream fitting.

\paragraph{Seed protocol.} In the high-data experiment, one seed value enters the stochastic components: the permutation that draws the validation split from the training pool,
each member's initialization and batch order, the half-split of the stacking
comparison, the subsamples used to tune the kernel stages, and the conformal
calibration draw; the pool/test boundary remains fixed. In the low-data protocol the first 1000 fitting and next 250 validation rows remain fixed across seeds. ClimSim provides no fixed test split,
so its seed also draws the $20000$-row test block out of the full set
(\texttt{campaign/climsim\_seeded.py}). ClimSim seed variation therefore includes test sampling, whereas the
structural-mechanics and OCO-2 test blocks are fixed. Tables report the
mean and standard deviation across seeds.

\paragraph{Calibration and readout controls.} The conformal comparison
uses the same calibration split for the constant-radius,
disagreement-scaled and $P_\lambda$-scaled bands. Each seed's
$P_\lambda$ is computed from its selected correction kernel
(Section~\ref{sec:uq}). The OCO-2 readout control repeats network training
at ten seeds per band and fits ridge and kernel readouts to the same
frozen features within each repetition (Table~\ref{tab:oco2ridge}).
The ClimSim kernel-budget comparison varies the number of fitting rows
under the protocol of Supplement~\ref{supp:oco2}.

\paragraph{Cost of individual kernel stages.} The parameter counts of
Table~\ref{tab:members} are the networks' weights alone. A deployed
structural-mechanics pipeline also stores the coefficient matrix of the
kernel member and that of the residual correction, each $19000\times1681$,
about $\costCoefM$ million numbers apiece ($255$\,MB each in double
precision), and evaluates a kernel row against the $19000$ training loads
for every query. At these dimensions, an eight-thread calculation on the same host uses
synthetic values to time the arithmetic in each operation. Block-filled
Gram construction takes $\costGramS$\,s, Cholesky factorization takes
$\costCholS$\,s, and triangular solves with $1681$ right-hand sides take
$\costSolveS$\,s. The peak resident memory is $\costPeakBlockGb$\,GiB against
$\costPeakNaiveGb$\,GiB after full-matrix construction. The lower peak uses block-filled construction; the
experimental correction implementation, \texttt{correct\_stack.py},
uses full-size temporaries.
At inference one kernel stage costs
$\costQueryOneMs$\,ms for a single query and $\costQueryBatchMs$\,ms per
query in batches of a thousand. These timings measure one kernel stage.
The full six-member pipeline evaluates two such stages plus five neural
members with reflection averaging; its total inference latency was not
measured by this diagnostic.

\paragraph{Experimental schedules.} The first structural-mechanics experiment
used ten seeds on four CPU instances. Kernel ridge, the metric-loss MLP, the
normalized-MSE MLP and the refiner completed their schedules at every seed.
A 36-hour per-task limit interrupted all ten FNO schedules. The reported
FNO predictions use the best observed validation checkpoint at each seed,
selected between epochs 10 and 70 of the scheduled 200. Validation error
increased after the selected checkpoint in every observed training curve;
the incomplete schedules do not establish performance after full training.
No UNet was trained in this experiment. Four- and five-member stacks, residual
corrections, second-moment measurements and conformal calibrations were
completed at every seed. Kernel-stage timings distinguish the standard and larger host classes
(Table~\ref{tab:compute}); Section~\ref{sec:cost} compares the memory
requirements of full-matrix and block-filled Gram construction.

The second experiment used the same ten seeds on one host with 20 physical
cores and 40 hardware threads. Repeated fits of kernel ridge and the three
MLP-family members agree with the first experiment's test errors to within $6\times10^{-8}$ in the relative-error unit
for the two MLPs and the kernel ridge and $6\times10^{-6}$ for the refiner.
The FNO and UNet each use a completed 100-epoch cosine schedule. Validation
selects the final checkpoint for nine FNO runs and all ten UNet runs, and
the epoch-79 checkpoint for the remaining FNO run. Each cosine schedule is normalized to 100 epochs and ends at $10^{-6}$. The six-member analysis includes per-pixel and global convex stacking,
residual correction, second moments and conformal calibration at every
seed. The paired comparisons of Section~\ref{sec:members-seeds} repeat
these analyses on the five members excluding the UNet; the half-split comparison selects global convex weights at one seed
of the five-member experiment, giving a corrected global convex stack. Table~\ref{tab:memcost} gives the training costs of each member.

The OCO-2 raw-input kernel size sweep uses ten seeds per band.
The separate O2 learning curve in Section~\ref{sec:scaling} uses one seed;
the ClimSim curves use five, three and three seeds at their respective
training sizes. The low-data pipeline uses ten seeds. The OCO-2 head
pipeline uses ten seeds on each of the three bands, giving thirty
band-seeds at a matched $250$-epoch budget, with one additional run at $750$
epochs to examine sensitivity to training duration. The combination rules, empirical floors and ensembles in
Section~\ref{sec:oco2} use member predictions and per-sample errors
from each band-seed. The readout control uses the same seeds and data-split protocol as
Table~\ref{tab:oco2}; small differences in weighted-network errors
reflect repeated training.
Seven of the thirty band-seeds were also repeated on different hardware:
six of the ten rows reproduce to the printed precision across the two runs,
including the kernel rows, the flat-metric network, its feature head and the
reported combination, while the weighted-loss networks
and their kernel heads differ by up to $0.39$ percentage points in reduced
error. Numerical reduction order can vary across hardware despite fixed
initialization and data order. These differences are smaller than the
corresponding between-seed standard deviations.

\paragraph{Prediction precision.} Errors recomputed from the stored member
predictions use the same metric and seed-specific validation partitions as
the original evaluations. On the fixed test block, all forty comparisons
with the original reflection-averaged errors agree to between $10^{-14}$
and $6\times10^{-8}$, consistent with the storage precision.

\paragraph{Residual data and rounding.} The residual and selection analyses
of Sections~\ref{sec:diversity} and~\ref{sec:members-seeds} use
$r_m=f_m-y$. For the first experiment, each of the ten seeds has one compressed NumPy archive,
\texttt{sm\_residuals\_s$k$.npz}, containing the test-block residual of each of the
five members and of the reported pipeline (per-pixel stack plus kernel
correction) as a $20000\times1681$ array, the validation-split residual of each
member as a $1000\times1681$ array, the target norms $\|y\|$ of every sample on
both splits, and the index arrays of both splits; the residuals are stored in
IEEE half precision (\texttt{float16}), the norms and indices at full width.
Across all 110 arrays, the mean validation and test errors computed
from the half-precision residuals differ from the original values by at most
$5.4\times10^{-8}$ in relative-error units. The largest stored residual
is $192$, compared with target norms of $3.7\times10^3$ to
$1.7\times10^4$. The measured rounding effect on these mean errors is far
below their seed spread. The archives occupy $384$\,MB per seed, $3.84$\,GB in all.
The archive manifest specifies the member, seed, split, dimensions,
precision, checksums and error summaries. The member errors in
Table~\ref{tab:members} are sample means of $\|r_m\|/\|y\|$. The
normalized validation residuals $r_m/\|y\|$ determine the second-moment
matrix $S$ and simplex weights; the test risks are computed from
$\|\sum_m w_m r_m\|/\|y\|$ using the sum-to-one identity. The same arrays
supply the dispersion factor in Proposition~\ref{prop:secmom}, pairwise
correlations, the equicorrelation summary, equal-weight and fitted
ensemble errors, spatial error maps, and the hard-sample comparison in
Section~\ref{sec:diversity}. Conformal scores use $\|r\|$ and the same
seed-specific calibration subset of the test block. All these quantities
inherit the half-precision rounding described above.
The per-pixel affine refit and kernel correction require the dataset and
member predictions on the relevant fitting splits. The kernel's posterior scale also requires its Gram matrix.
The five-member archive
does not contain the second experiment's completed FNO or added UNet.

\paragraph{Runtime.} Table~\ref{tab:compute} reports wall times from
the first experiment. Each task used six threads, with five tasks sharing an
instance. Results are separated by host class: seeds 2, 4, 8 and 9 used the
larger instance for the kernel stages. A factorization, solve and validation
prediction at $n=19000$ took $20$--$35$\,s there, compared with
$42$--$113$\,s on the standard instances. The FNO's incomplete schedules
are summarized by seconds per epoch from cumulative ten-epoch timers.
Uninterrupted intervals took $467$--$625$\,s per epoch on standard
instances and $348$--$410$\,s on the larger instance. Intervals affected by
shared-host contention or suspension reached $10^4$\,s per epoch; the table
therefore reports both minimum and median rates. Selected checkpoints were
reached after $1.3$--$8.0$\,h at five seeds with uninterrupted timers, and
$9.0$--$16.9$\,h at the remaining five. Full-test inference latency and
peak memory were not recorded. Per-seed timings are provided in \texttt{runs/timing\_harvest.json}.

\begin{table}[t]
\centering
\caption{Wall times for the first structural-mechanics experiment
(all stages on CPU, six threads per task,
five tasks sharing each instance). Network rows: the trainer's elapsed
minutes from the first epoch to the final evaluation, as the range over seeds,
with the scheduled epochs and the selected checkpoint epoch (range over the
ten seeds). The FNO row gives seconds per epoch from ten-epoch timing intervals, minimum and median over all intervals (count in parentheses), since
no seed completed the schedule. Kernel rows: the complete stage (fifteen grid
cells and the refit) and its individual factorization-and-solve operations, as ranges,
with the number of timed cells or folds in parentheses. The two columns are
the two instance classes of the experiment, with the number of seeds in parentheses.}
\label{tab:compute}
\small
\setlength{\tabcolsep}{3.5pt}
\resizebox{\textwidth}{!}{%
\begin{adjustbox}{max width=\linewidth}
\begin{tabular}{llcc}
\toprule
Stage & Epochs: scheduled / selected & Standard instances (6 seeds) & Larger instance (4 seeds) \\
\midrule
Metric-loss MLP & 400 / 39--79 & 129--158 min & 26--35 min \\
MSE MLP & 120 / 99--119 & 9.7--11.9 min & 6.6--11.1 min \\
Refiner & 100 / 99 & 9.9--13.1 min & 7.9--33.7 min \\
FNO, per epoch & 200 / 10--70 & 467 s min, 581 s median (88) & 348 s min, 437 s median (56) \\
Kernel ridge, grid and refit & --- & 17.0--26.0 min & 6.1--7.4 min \\
\ \ one factor-and-solve, $n=19000$ & --- & 42--113 s (90) & 20--35 s (60) \\
\ \ out-of-fold fold, $n=14250$ & --- & 32--66 s (24) & 13--22 s (16) \\
\bottomrule
\end{tabular}
\end{adjustbox}}
\end{table}

\paragraph{Member costs.} Table~\ref{tab:memcost} gives training costs
for the second experiment over ten seeds, with all schedules completed
on one host. Thread counts differ across members because the
host was shared with the rest of the experiment. We report thread-minutes,
elapsed wall time multiplied by the configured thread count. This is allocated parallelism; CPU utilization was not measured, and the timings include contention on the shared host. The FNO uses about $7900$ thread-minutes
against $54$ for the MSE MLP while improving single-member error by about
$0.03$ percentage points. The UNet uses about $115$ times the MSE MLP's
thread-minutes and has higher single-member error, although it improves the
ensemble. Parameter counts alone do not explain these differences: the FNO
has 6.45M parameters, the refiner 6.64M, and the UNet 4.38M. Convolution over
the full field and dense prediction from the load vector have different
computational costs; the table measures their implementations at the stated schedules.

\begin{table}[t]
\centering
\caption{Training costs of the second experiment on a shared CPU host with 20 physical cores and 40 hardware threads,
every member trained to the end of its schedule; ranges over the ten seeds.
Wall minutes are elapsed training times; thread-minutes is wall clock
times the thread count (10 for the convolutional members, 10 to 16 for the
MLPs); the kernel stage does not record a
thread count. Test errors use reflection averaging for the neural members,
as in Table~\ref{tab:members}. Per-seed values are provided in
\texttt{campaign/collected/dgx/runtime\_ledger.csv}.}
\label{tab:memcost}
\small
\setlength{\tabcolsep}{4.5pt}
\begin{adjustbox}{max width=\linewidth}
\begin{tabular}{lccccc}
\toprule
Member & Parameters & Epochs & Wall (min) & Thread-min & Test error (\%) \\
\midrule
Kernel ridge, grid and refit & --- & --- & 5.8--10.0 & --- & 5.193--5.200 \\
Metric-loss MLP & 4.91M & 400 & 8.8--19.2 & 132--192 & 4.809--4.858 \\
MSE MLP & 4.91M & 120 & 2.6--5.8 & 39--58 & 4.703--4.719 \\
Refiner & 6.64M & 100 & 2.6--6.1 & 41--61 & 4.721--4.742 \\
FNO & 6.45M & 100 & 662--836 & 6618--8355 & 4.665--4.697 \\
UNet & 4.38M & 100 & 573--689 & 5734--6892 & 4.683--4.793 \\
\bottomrule
\end{tabular}
\end{adjustbox}
\end{table}

\paragraph{Stacking.} The five-member fitted-weight baseline uses
least squares on flattened full-validation predictions, followed by clipping
negative coefficients and normalizing the weights to sum to one
(\texttt{campaign/ensemble\_seeded5.py}). The six-member fitted-weight
baseline uses full-validation mean relative error and softmax weights:
\texttt{stack\_correct.py} starts from uniform weights and takes 2000
Gaussian proposals in logit space, with initial step $0.5$, decay factor
$0.999$ and random seed 0. The separate global-weight fit used to judge
per-pixel stacking uses half the validation set, 600 proposals, initial
step $0.3$ and decay factor $0.997$. The
simplex minimizer of the second-moment matrix $S$
(Proposition~\ref{prop:secmom}) is used only in the pool analyses.
Per-pixel stacking fits an affine combination of member predictions at
each grid point, with ridge parameter $10^{-3}$. The weights are fitted
on half the validation split and selected if their error on the other
half is lower than that of global weights. The selected stack is then
refitted on the full split.

\paragraph{Kernel stages.} The kernels are Mat\'ern-$5/2$ functions of
standardized inputs. The residual
correction searches the scale grid
$\{1,2,4\}\times$ the median pairwise distance (estimated on 2000
fitting rows) and the nugget grid $\{10^{-6},10^{-5},10^{-3}\}$ (scaled by $n$); the
low-data chain of Table~\ref{tab:lowdata} uses $\{0.5,1,2,4\}$ with nuggets
$\{10^{-7},10^{-5},10^{-3}\}$, with the median computed on all 1000 fitting
rows. Candidate fits use $\min(8000,n)$ fitting rows, where $n$ is the
fitting-set size: 8000 rows in the high-data protocol and all 1000 rows in
the low-data protocol. Hyperparameters are selected on validation, then
refitted on all fitting rows in double precision. The pure kernel
baseline (Table~\ref{tab:main}) uses the grid
$\{0.5,0.75,1,1.5,2\}\times$median and nuggets $\{10^{-8},10^{-6},10^{-4}\}$
directly at $n=19000$. The feature-space stage concatenates and
standardizes penultimate activations, then uses the same kernel family
and tuning protocol. These features are spatially averaged pre-projection
channels for the FNO (128 coordinates), mean-pooled encoder tokens for the
transformer (192), and the residual-trunk output for the MLP (1024).

\paragraph{Low-data correction sequence.}
The low-data sequence uses one trained MLP as its initial mean. The candidate predictors are the neural mean,
the mean plus an input-space correction, that predictor plus a
feature-space correction, and the resulting predictor plus a second
input-space correction. Each correction is fitted to the remaining
training residuals. The final predictor is the candidate with the lowest
error on all 250 validation rows.

\paragraph{Low-data multiscale kernel.}\label{supp:multiscale-lowdata}
The kernel baseline in Table~\ref{tab:lowdata} uses the additive residual
construction in \texttt{campaign/krr\_multiscale\_lowdata.py}. Inputs are
standardized by the coordinatewise mean and standard deviation of the
$n=1000$ fitting rows. Initialize $g_0(x)=\bar y$, the fitting-target mean.
At stage $t$, each candidate fits the current residual matrix
$R_t=Y-g_{t-1}(X)$ and proposes
\[
 g_t(x)=g_{t-1}(x)
   +k_{\ell_t}(x,X)(K_{\ell_t}+n\lambda_t I)^{-1}R_t.
\]
The Mat\'ern-$5/2$ length scale is selected from
$\{0.25,0.5,1,2,3\}d_{\rm med}$, with
\[
d_{\rm med}=\left(\operatorname{median}_{i<j}\|x_i-x_j\|_2^2\right)^{1/2}
\]
computed on the standardized fitting inputs. The nugget grid is
$\{10^{-8},10^{-6},10^{-4}\}$. At each stage the candidate with the lowest
mean relative-$L^2$ error on the fixed 250-row validation split is selected.
Its increment is retained with coefficient one only if that error improves
by more than $10^{-5}$ in relative-error units. Otherwise the procedure
stops. At most four stages are retained, and all fits use the same 1000
fitting rows. The procedure is deterministic for this fixed split.

\paragraph{OCO-2 network and feature heads.}\label{supp:oco-network}
The implementation in \texttt{campaign/jpl\_seeded.py} uses a width-384
residual MLP. A linear input map and SiLU activation precede three
residual layers $h\mapsto h+\operatorname{SiLU}(Wh+b)$ and a linear output
map to the 40 standardized PCA coordinates. The final 384-dimensional
hidden vector supplies the frozen features for the kernel and ridge
readouts. Training uses AdamW with learning rate $10^{-3}$, weight decay
$10^{-5}$, batch size 512 and a cosine schedule to $10^{-6}$ over 250
epochs. Validation is evaluated every 25 epochs and at the final epoch;
the checkpoint minimizing the corresponding training metric on validation
is retained. The three losses are mean relative error in standardized PCA
coordinates, mean relative error after diagonal scaling by the PCA
coefficient scales, and mean relative error after the retained-PCA radiance
reconstruction. The additional 750-epoch run uses the same architecture and
optimizer, with the cosine schedule extended to that budget.

For each OCO-2 kernel head, features are standardized by fitting-set
moments. The Mat\'ern-$5/2$ scale grid is
$\{0.5,1,2,4\}d_{\rm med}$ and the nugget grid is
$\{10^{-8},10^{-6},10^{-4}\}$. A seeded 6000-row fitting subsample is used
both to estimate $d_{\rm med}$ and to fit candidates, which are selected on
the full validation split using the associated loss. The selected pair is
refitted on all 18000 fitting rows. Raw-input kernels use the same grid,
subsample size and refitting protocol. The ridge control includes an
intercept and selects $\lambda$ from
$\{10^{-10},10^{-8},10^{-6},10^{-4},10^{-2},1\}$ on validation, with
normal equations regularized by $n\lambda I$.

\paragraph{OCO-2 input sensitivity.}\label{supp:oco-input}
Let $X_s$ contain the fitting inputs standardized by their coordinatewise
fitting means and standard deviations, and let $Z$ contain the released
standardized PCA targets. The sensitivity estimator in
\texttt{campaign/jpl\_seeded.py} is the least-squares matrix
\[
 A=X_s^{\dagger}Z,\qquad
 a_j=\|A_{j,:}\|_2,\qquad
 w_j=\frac{a_j}{d^{-1}\sum_{k=1}^{d}a_k}.
\]
The adapted kernel uses distances between $\operatorname{diag}(w)x_s$.
The weights are applied after standardization, so subsequent input
normalization does not remove them. The estimator is refitted using only
the available fitting rows at each training size. The rank summaries in
\texttt{campaign/scaling\_seeded.py} instead use
$B=X_{\rm tr}^{\dagger}Z$ on all 18000 fitting rows, where $X_{\rm tr}$
contains the released input coordinates without the additional fitting-set
centering and standardization. For the singular values $\sigma_j$ of $B$,
the 99\% count is the smallest $r$ satisfying
$\sum_{j\le r}\sigma_j^2\ge0.99\sum_j\sigma_j^2$, and the reported
participation ratio is $(\sum_j\sigma_j)^2/\sum_j\sigma_j^2$.

\paragraph{Separate O2 learned-metric experiment.}\label{supp:oco-learned-metric}
The experiment in \path{learnable_kernel_gpu.py} compares an isotropic
metric, a learned diagonal metric, and
$M=L^\top L+\operatorname{diag}(\exp(2b))$, with $L\in\mathbb R^{6\times d}$.
The diagonal model has $d$ metric parameters and the low-rank model $7d$.
After a fixed seed-0 permutation, kernel fitting uses 6000 pool rows and
validation uses the next 2000. Metric learning uses 400 Adam steps at
learning rate $0.05$, with 256 fitting rows sampled with replacement per
step. For a batch $b$ and its first half $c$, let $K_b$ and $K_c$ be the
Mat\'ern Gram matrices and let $Z_b$ and $Z_c$ contain their reduced
targets. The regularized kernel-flows objective is
\[
 \rho=\max\left\{10^{-6},\,
 1-\frac{\operatorname{tr}\!\left[Z_c^\top(K_c+\eta I)^{-1}Z_c\right]}
 {\operatorname{tr}\!\left[Z_b^\top(K_b+\eta I)^{-1}Z_b\right]+10^{-30}}
 \right\},\qquad \eta=\exp(\delta)+10^{-8}.
\]
This is the regularized half-batch criterion of
\citet{owhadi2019kernel}, with the quadratic forms summed over outputs.
The parameter $\delta$ is learned jointly with the metric.
With the learned metric fixed, validation
selects a multiplier of its feature map from
$\{0.1,\allowbreak 0.25,\allowbreak 0.5,\allowbreak 1,\allowbreak 2,\allowbreak 4,\allowbreak 8\}$ and a nugget from
$\{10^{-8},10^{-6},10^{-4},10^{-2}\}$. The final Mat\'ern-$5/2$ solve uses
the same 6000 fitting rows.

The raw-state experiment has $d=20$, so the low-rank term adds 120
parameters. The feature experiment uses a separate width-256 MLP with two
residual SiLU blocks, trained by standardized-target MSE for 600 epochs
on the 20000-row pool before the kernel split. Its $d=256$ feature map gives
1536 additional low-rank parameters. These experiments use different
training budgets and representations from Table~\ref{tab:oco2}; in the
feature experiment, the kernel-validation targets have already contributed
to neural training. They provide a separate comparison of metric choices
within those fitted representations.

\paragraph{Uncertainty.} The posterior standard deviation in
\eqref{eq:power} is computed by triangular solves using the Cholesky
factor of $K+n\lambda I$.
Conformal calibration uses a random 1000-sample subset of the public test block and evaluation uses the disjoint remainder. Model-fitting optimizers use training and validation. The marginal coverage guarantee of Proposition~\ref{prop:conf} requires the predictor and scale to be fixed independently of exchangeable calibration and evaluation observations. The exact Beta reference law additionally assumes i.i.d. continuous scores. The miscoverage levels $\alpha=0.10$ and $0.05$ have $k\le m$; for smaller miscoverage with $k=m+1$, the mathematical quantile is $+\infty$, not the largest finite calibration score.
